\section{Automatic scoring from shared construction logic}
\label{sec:scoring}

Each item model writes two things at once: the prompt shown to the language model and the checklist used to score the diagram (\cref{fig:cogen}). Because the finished diagram is a set of coordinates, the engine can decide each check directly. ``Is the angle at $H$ a right angle?'' is settled by a single calculation on the coordinates. No judge is asked whether the picture appears square. This makes scoring inspectable and inexpensive. It does not make scoring complete. The checklist's validity still depends on whether it captures everything the prompt asks for and the visual conventions a reader expects. We examine both limits below.

\paragraph{The question and its answer key are written together.}
Each item and its answer key come from the \emph{same} short program. Because the two are emitted together, they are far less likely to drift apart through manual editing. Errors in the item model itself remain possible. An author can still omit a predicate, encode the wrong one, or write wording whose intended constraint the checklist does not fully capture. Re-running the item model with a new random seed produces a fresh, parallel item. The checklist is never shown to the language model while it works.

\begin{figure*}[t]
\centering
\begin{tikzpicture}[
    every node/.style={font=\small},
    func/.style={draw, rounded corners, fill=blue!8, align=center,
                 text width=70mm, minimum height=10mm, font=\ttfamily\small},
    outbox/.style={draw, rounded corners, fill=green!6, align=center,
                text width=53mm, minimum height=18mm, inner sep=4pt},
    role/.style={font=\scriptsize\itshape, color=gray!50!black, align=center},
    >=Stealth, line width=0.5pt
]
  \node[func] (fn) {altitude\_template($A,B,C,H$,\ ang$_A$,\ ang$_B$):};
  \node[outbox, anchor=north] (prompt) at ($(fn.south)+(-4.0,-1.4)$) {%
    \textbf{Prompt (shown to the language model)}\\[2pt]
    \scriptsize ``Draw an acute triangle $ABC$ with $\angle A = 60^{\circ}$, $\angle B = 70^{\circ}$.
    Construct the altitude from $C$ meeting $AB$ at $H$.''};
  \node[outbox, anchor=north] (rubric) at ($(fn.south)+(4.0,-1.4)$) {%
    \textbf{Answer key (used to score)}\\[2pt]
    \scriptsize\ttfamily right\_angle($C, H, A$)\\
    point\_on\_segment($H, A, B$)};
  \draw[->] (fn.south) .. controls +(0,-0.6) and +(0,0.6) .. (prompt.north);
  \draw[->] (fn.south) .. controls +(0,-0.6) and +(0,0.6) .. (rubric.north);
  \node[role, below=1mm of prompt, text width=53mm] {what the language model sees};
  \node[role, below=1mm of rubric, text width=53mm] {hidden from the language model};
\end{tikzpicture}
\caption{One item model writes both the question and its answer key. The two cannot drift apart through separate editing. An item model can still leave a check out. This key confirms the altitude but not the two stated angles. That gap is the main finding of the human study reported in this section.}
\label{fig:cogen}
\end{figure*}

\paragraph{A small, readable checklist.}
Each encoded requirement uses a small, fixed vocabulary of checks (\cref{tab:predicates}): right angle, midpoint, parallel lines, point on segment, and so on. There are 15 such checks in the current version. The construction is computed symbolically; each check is then decided numerically on the coordinates, within a small tolerance (\cref{sec:benchmark}). A diagram passes only if every check on its list passes. The list defines what the automatic scorer evaluates. It plays the role a specification plays in formal methods: a statement of what must hold, written so that a machine can check it. The check here is numeric and applied to one figure at a time, not a proof over all figures. For that reason we say \emph{verifiable} rather than \emph{verified}.

\begin{table}[t]
\centering
\footnotesize
\setlength{\tabcolsep}{4pt}
\begin{tabular}{@{}l p{0.5\columnwidth}@{}}
\toprule
\textbf{Check} & \textbf{What it confirms}\\
\midrule
\propname{right\_angle}                     & an angle is a right angle\\
\propname{midpoint}                         & a point is halfway along a segment\\
\propname{collinear} / \propname{parallel}  & points lie on a line / lines never meet\\
\propname{point\_on\_segment}               & a point lies between two endpoints\\
\propname{equal\_lengths}                   & two segments have equal length\\
\propname{angle\_equal}                     & an angle equals a stated value\\
\bottomrule
\end{tabular}
\caption{A sample of the 15 checks in the scoring vocabulary. Each is decided on the computed coordinates. An item's list of checks is what a passing answer must satisfy.}
\label{tab:predicates}
\end{table}

\paragraph{Does the automatic score match a human's?}
An automatic score is useful only if it agrees with a careful human grader, so we tested that agreement. We fixed a few blind spots found on an early pilot. We then ran a pre-registered study on a later run of the benchmark.\footnote{Protocol: \texttt{docs/human\_study\_protocol.md} in the code repository, \url{https://github.com/RenaissancePhilanthropy/geometry-diagram-generator}.} We set the bar in advance: the automatic score had to match the human verdict at Cohen's $\kappa \ge 0.70$, on a scale where $1.0$ is perfect agreement and $0$ is chance. Two graduate-level graders then scored 200 diagrams by hand, without seeing the automatic verdict. The 200 were drawn from that run stratified by automatic verdict: every uncertain and every rendered failing diagram, plus passes balanced across models, modes, and difficulty tiers. Each grader rated independently, and disagreements were then discussed to a single consensus verdict per diagram. We compare the automatic verdict with that consensus, both collapsed to pass or fail, with a partial human verdict counted as a fail.

We did not clear that bar. Agreement was moderate, Cohen's $\kappa = 0.46$, well short of the pre-registered $0.70$. We treat this as a central limitation. The disagreements run \emph{one way}. In 36 of the 200 cases the automatic scorer passed a diagram the graders failed. No automatic failure received a passing human verdict. That is the more serious direction once a score is read as a certificate of correctness, because every disagreement is a false positive. We therefore state it plainly: \textbf{the pass rates in \cref{sec:benchmark} may overestimate correctness}. Inspecting the 36 disagreements traced them to the checklist rather than to the ratings. They concentrate in items with an added construction, like our altitude, whose checklist confirms the construction but does not re-check the prompt's stated angles. The two graders agreed with \emph{each other} at the $0.70$ we had asked the scorer to reach. The fix is to add the missing numeric checks to those item models. That revision was not applied to the results reported here. Instead, \cref{sec:benchmark} reports, alongside each automatic pass rate, the share of automatic passes the graders confirmed, as a sensitivity calculation.

\paragraph{What the checks cannot see.}
That particular gap closes with more checks. A deeper one does not. The checks read the geometry, the coordinates, and nothing else. They never look at how the figure is drawn. A diagram can satisfy every requirement and still mislead a reader. A right-angle mark can be placed outside the corner. An angle can be labelled on the wrong side. A figure can pass every check and still be laid out so poorly that a reader cannot follow it. This is the converse of the language model-based judge in \cref{sec:motivation}. A judge sees how a picture looks but cannot check its geometry. Our checks test the geometry but cannot see how it looks. The two are complementary. The checks say whether the computed geometry is correct. Whether the drawing presents that geometry faithfully and clearly still requires visual review.
